\section{LITERATURE REVIEW}
\label{ch:literature}

\subsection{Theoretical Background}
\label{sec:theory}

An automated post-harvest grading and sorting machine sits across three fields: the ripening physiology that determines what is being graded, the detection architecture that performs the grading, and the control design that converts a grading decision into a physical action. This section reviews each in turn.

\subsubsection{Ripening Physiology and Cultivar-Specific Grading}
\label{sec:ripening}

Ripening in tomato is climacteric. Under the influence of endogenous ethylene the exocarp loses chlorophyll while carotenoid pigments, principally lycopene, accumulate, so that surface colour moves from uniform green to deep red; in parallel, cell-wall pectins are progressively degraded and the fruit softens \citep{abekoon2024}. Softening is the point at which the fruit becomes commercially fragile, because reduced firmness raises susceptibility to bruising and to the surface damage through which decay organisms enter.
% TODO: add one dedicated post-harvest physiology reference here (a text or
% review covering pectin methylesterase and polygalacturonase). The Word
% draft cited allo2025 for this, which is a CNN classification paper and
% does not support the claim.

Because surface colour tracks this internal progression, it is the basis of standard visual grading. The \acrfull{USDA} standard defines six stages: green; breaker, at not more than 10\,\% red surface; turning, at 10--30\,\%; pink, at 30--60\,\%; light red, at 60--90\,\%; and red, above 90\,\% \citep{usda1991}. These stages are defined for grading and reporting, not for machine vision, and the two narrowest bands sit exactly where a camera is least reliable: under packhouse lighting, belt motion and specular reflection from the glossy cuticle, distinguishing 8\,\% red surface from 12\,\% is not a stable measurement.

Generic standards also travel poorly between cultivars. \citet{nalupano2024} retrained a MobileNetV2 network specifically for the Philippine Kinalabasa variety, reaching 94\,\% accuracy but only 64\,\% specificity, which illustrates both that cultivar-specific training is necessary and that headline accuracy can conceal weak class separation. For Sri Lanka, \citet{abekoon2024} studied the Padma variety directly. One hundred greenhouse-grown fruits, twenty-five from each of four districts, were held in a facility reproducing the mean conditions measured across sixty Sri Lankan marketplaces (29\,\textdegree C, 80\,\% relative humidity, 810 lux) and photographed daily until deterioration, yielding 12,342 images. Their reported stage durations are reproduced in Table~\ref{tab:abekoon_stages}. That work establishes both a cultivar-specific colour-to-time relationship and the six-stage frame within which the present study defines its own classes.

\begin{table}[H]
\caption{Post-harvest stage durations reported for the Padma variety under simulated market conditions}
\label{tab:abekoon_stages}
\centering
\begin{tabular}{lcc}
\toprule
\textbf{Maturity stage} & \textbf{Days in stage} & \textbf{Post-harvest day range} \\
\midrule
Green       & 2  & 0--2   \\
Breakers    & 2  & 3--4   \\
Turning     & 3  & 5--7   \\
Pink        & 4  & 8--11  \\
Light red   & 8  & 12--19 \\
Red         & 12 & 20--31 \\
\midrule
\textbf{Total} & \textbf{31} & \\
\bottomrule
\end{tabular}

\vspace{2mm}
{\footnotesize Source: \citet{abekoon2024}, Table 4. Conditions: 29\,\textdegree C, 80\,\% relative humidity, 810 lux; 100 fruits.}
\end{table}

\subsubsection{Real-Time Single-Stage Object Detection}
\label{sec:detection}

Early agricultural vision systems used hand-crafted colour and texture descriptors with shallow classifiers. \citet{santoso2024} compared \acrshort{RGB}, \acrshort{HSV} and CMYK representations using a correlation-coefficient matching method, finding CMYK the most discriminative with a correlation of 0.87. \citet{liu2019} combined histogram-of-oriented-gradients descriptors with a \acrfull{SVM} and a colour-based false-positive removal stage. \citet{garcia2019} classified ripening stage from pixel colour statistics. These methods are computationally light, but they depend on descriptors chosen in advance and degrade when illumination, shadow or background departs from the conditions in which they were tuned \citep{ifmalinda2023}.

Convolutional networks removed the need to choose descriptors by learning them from data. \citet{allo2025} trained a convolutional network on 3,600 public images across three ripeness classes, and \citet{centino2025} obtained 95.52\,\% on a binary ripe/unripe task using background removal and colour-space conversion before classification. \citet{waseem2025} pursued the same task under a computational constraint, using a ResNet-18 backbone with pruning and quantisation to keep maturity classification tractable on limited hardware. All three, however, classify a whole image. A sorting line needs localisation as well, because a frame may contain more than one fruit and each requires its own decision.

Object detectors supply that. Two-stage detectors such as Faster \acrshort{RCNN} first propose regions using a \acrfull{RPN} and then classify them, which is accurate but too slow for a moving belt. Single-stage detectors of the \acrshort{YOLO} family instead regress bounding boxes and class probabilities in one forward pass \citep{redmon2016}. YOLOv8 is anchor-free and uses a cross-stage partial module with two convolutions for multi-scale feature fusion, giving the inference speed required for real-time operation \citep{terven2023}.

Tomato-specific work concentrates on adapting this family to difficult imaging conditions. \citet{yang2025} and \citet{fan2026} modified YOLOv8 for occlusion and lighting variation in greenhouse and orchard scenes; \citet{huang2025} added a small-target detection head and multi-scale fusion to YOLOv10; \citet{wang2026} and \citet{ding2026} produced lightweight variants for deployment on constrained devices. \citet{safira2026} compared YOLOv8 directly against a real-time detection transformer for tomato ripeness and provides the most direct published basis for preferring a \acrshort{YOLO} detector for this task. Almost all of this work targets harvesting robots operating in the field or greenhouse, where the dominant problems are occlusion and cluttered backgrounds. A packhouse belt presents the opposite situation: a controlled background and a single fruit in view, but a decision that must be reached within the seconds the fruit is visible and then executed mechanically.

\subsubsection{Learned Imaging Conditions as a Confounding Cue}
\label{sec:domainshift}

A detector learns whatever feature separates its classes in the training data, including features that have nothing to do with the object. If the images of one class were captured under different lighting or against a different background from the rest, the network can learn that difference instead of the intended visual property, and will appear accurate on a test split drawn from the same distribution while failing on live input. The risk is highest when one class is scarce and is topped up from an external source, which is exactly the situation for a defect class. This is a well-recognised hazard, and the standard mitigations are to hold capture conditions constant across all classes and to increase minority-class representation by augmentation rather than by importing outside images \citep{shorten2019, buslaev2020}. It is nevertheless rarely reported in applied agricultural vision papers, which tend to publish final metrics rather than the diagnostic history behind them. Section~\ref{sec:annotation_qc} documents an instance encountered in the present study.

\subsubsection{Conveyor Actuation and Control}
\label{sec:actuation}

Converting a classification into a physical diversion requires the gate to open at the moment the fruit reaches it. The straightforward approach computes the delay from belt kinematics. For gate $i$ at distance $d_i$ from the camera, with belt velocity $v$, inference latency $t_{\text{inf}}$ and transmission latency $t_{\text{tx}}$, the delay before actuation is given in Equation~\eqref{eq:timed}.

\begin{equation}
    t_{\text{act},i} = \frac{d_i}{v} - t_{\text{inf}} - t_{\text{tx}}
    \label{eq:timed}
\end{equation}

Equation~\eqref{eq:timed} is elementary kinematics, and its weakness is the assumption that $v$ is constant and known. On a real rig, belt tension, load, drive slip and motor torque all vary, and because $d_i$ is the full camera-to-gate distance the resulting position error grows with distance down the line, affecting the last gate most. The alternative is to detect the fruit's arrival rather than predict it, by placing a presence sensor at each gate and holding pending classifications in a queue until the corresponding sensor is triggered. Existing sorting rigs use belt sensors, though generally for a different purpose: in the system of \citet{moya2025}, sensors positioned along the conveyor detect fruit and trigger image capture, while the sorting servos act on the classification result. \citet{geetha2025} likewise couple detection to a servo-driven sorting prototype. Using a sensor at each gate to release that gate, so that the elapsed-time term is eliminated from the actuation path entirely, is the specific design adopted here.

\subsubsection{Embedded Integration and Wireless Command Transfer}
\label{sec:embedded}

Vision-based agricultural systems are commonly split between a host that performs inference and a microcontroller that drives actuators. \citet{patria2025} used an ESP32-CAM with an artificial neural network on a mobile field robot, and \citet{saxena2025} paired ESP32 sensor nodes with a Raspberry Pi running YOLOv8, using quantisation and pruning to gain a 35\,\% inference speed-up. Both illustrate the division of labour, though in both cases the microcontroller senses rather than actuates a sorting mechanism.

Where the link between host and controller is wireless, command delivery cannot be assumed. A serial-over-Bluetooth link on an operating rig is subject to intermittent loss, and in a sorting application a lost command is not merely a retransmission: it removes a fruit from the queue and misaligns every subsequent gate assignment. A command protocol for this purpose therefore requires per-command sequence numbering, explicit acknowledgement, and duplicate suppression so that a retransmitted command is acknowledged again but not acted on twice.
% NOTE: the Word draft attributed all of the above to patria2025. That paper
% is a field mobile robot with no conveyor, no gate and no queue, and does
% not support these claims. The argument is retained here on engineering
% grounds and is demonstrated in Chapter 3, not asserted from a citation.

\subsubsection{Shelf-Life Estimation and Operational Indicators}
\label{sec:shelflife_lit}

Shelf life is the period over which produce retains acceptable safety, firmness and sensory quality under defined storage conditions \citep{tarlak2023}. Two approaches to estimating it appear in the literature. The first measures it directly, by observing a cohort of fruit under controlled conditions and recording how long each stage lasts, as \citet{abekoon2024} did for the Padma variety. The second learns it, as \citet{geetha2025} did using deep models coupled to a sorting prototype. Non-destructive instrumental alternatives also exist: \citet{borba2021} used portable \acrfull{NIR} spectroscopy to assess tomato quality in the field, and \citet{liu2024} fused multiple sensing modalities for maturity assessment. These give internal quality measures that colour imaging cannot, at the cost of instrumentation that a low-cost packhouse rig cannot carry.

Whichever route is taken, the estimate becomes operationally useful only when it is aggregated. A packhouse manager needs throughput, class distribution, defect ratio and mean remaining shelf life across a batch, not a value for one fruit \citep{geetha2025}.

\subsection{Gaps in Knowledge}
\label{sec:gaps}

Three gaps follow from the review above. They are stated narrowly, because adjacent work exists in each case.

\textbf{No conveyor-integrated detection system exists for the Padma variety.} \citet{abekoon2024} characterised Padma ripening and post-harvest duration thoroughly, but the work is offline classification of static images of single fruits, with no detector, no localisation and no physical sorting. Conversely, conveyor-integrated systems such as \citet{moya2025} are calibrated for other cultivars in other production contexts. No published work reports a detection model trained on Padma fruit and evaluated on an operating sorting line, and no public Padma detection dataset spanning four ripening stages and a defect class exists to build one from.

\textbf{Gate release is generally predicted rather than detected.} Belt sensors are already used in sorting rigs, but predominantly to trigger image capture \citep{moya2025}. Actuation itself is typically timed from belt kinematics, per Equation~\eqref{eq:timed}, which makes sorting accuracy dependent on belt speed remaining constant and degrades progressively along the line. A control architecture in which each gate is released by its own presence sensor, with classifications held in a queue until arrival, removes that dependency but is not established in the tomato sorting literature.

\textbf{Grade-to-shelf-life mapping and line-level reporting are rarely combined with an operating rig.} This gap is the narrowest of the three, and must be stated carefully. \citet{geetha2025} do link deep-learning grading to shelf-life prediction and a servo-driven sorting prototype, so it cannot be claimed that no such system exists. What remains absent is a system in which the shelf-life values are measured for the specific cultivar being graded, rather than learned from generic data, and in which each classification and its associated remaining shelf life are logged to a persistent record that drives live batch-level operating indicators. It is that combination, for this cultivar, that the present study addresses.
